\section{A dynamic model with a uniform memory law}\label{sec:refreshing}
\subsection{The experiment and its resource convention}
Fix $d\ge1$, let $E_d=\{0,1,\ldots,d\}$, and write
$\nu_i=1/(d+1)$. The known parameters satisfy
\begin{equation}\label{eq:refresh-range}
 4/5\le\gamma\le1,\qquad 0\le\epsilon\le1/4.
\end{equation}
Prepare $X_0\sim\nu$. At each acquisition step the hidden state is updated
by the transition matrix
\begin{equation}\label{eq:refresh-transition}
 P_\gamma(i,j)=(1-\gamma)\1_{\{i=j\}}+\gamma\nu_j.
\end{equation}
A command $u_t\in[-1,1]^d$, chosen before the report, gives a binary report
$Y_t\in\{-1,1\}$ with likelihoods
\begin{equation}\label{eq:refresh-likelihood}
 \ell_y(u,i)=\begin{cases}
 1/2,&i=0,\\
 (1+y\epsilon\tanh u_i)/2,&1\le i\le d.
 \end{cases}
\end{equation}
Conditional on the hidden path and commands, the report randomizations are
independent. The resulting joint transition/report kernel is
$P_\gamma(i,j)\ell_y(u,j)$. It is an instance of
Definition~\ref{def:experiment}; all report probabilities lie between
$3/8$ and $5/8$.

For the acquisition lower bound the commands are independent uniform on
$[-1,1]^d$, independent of the hidden process and coding randomness. The
upper-bound filter will work on every admitted command/report sequence,
including adaptively chosen commands. It does not change the fixed
exploration law used for the lower bound.

At any selected time $t\ge1$ the acquisition can stop and one terminal
probe is executed without an intervening refresh. After the retained label
is formed, draw $J$ uniformly from $\{1,\ldots,d\}$. Probe $j$ produces
$Z\in\{0,1\}$ with
\begin{equation}\label{eq:refresh-probe}
 \Pp(Z=1\mid X_t=i,J=j)=1/2+\tau\1_{\{i=j\}},
 \qquad \tau=1/4.
\end{equation}
This is one noisy binary readout, not observation of $X_t$ or of the full
posterior. Its choice is independent of the acquisition. Each acquisition
and the terminal probe costs one recorded experimental step. Evaluating
maximum-checkpoint risk means that a run executes only its selected probe;
it does not execute counterfactual probes at all times.

Let $p_t=\Law(X_t\mid H_t)$, and write $p_{t,+}=(p_{t,1},\ldots,p_{t,d})$.
The probe's full-history success probability is $1/2+\tau p_{t,j}$.
If a decoder uses candidate coordinates $c_j$, its excess loss is
\begin{equation}\label{eq:refresh-distortion}
 \frac{\tau^2}{d}\norm{p_{t,+}-c}_2^2.
\end{equation}
Arbitrary probe predictions are allowed; projecting them onto the interval
$[1/2,1/2+\tau]$ cannot worsen loss. Thus the decoder coordinate formulation
loses no candidates for the infimum.

Let $R_{M,t}^{\rm av}$ and $R_{M,t}^{\rm max}$ be optimal checkpoint risks
under the stated exploration and over all admitted histories, respectively.
Let $\mathfrak R_{M,N}^{\rm av}$ and $\mathfrak R_{M,N}^{\rm max}$ be the
corresponding best maximum risks for one causal $M$-label encoder at times
$1,\ldots,N$. The clock, $d,\gamma,\epsilon$, and the finite read-only
program are known; no history or unknown calibration is stored in those
constants. At time zero one fixed initial state is allowed. For $N=\infty$
replace maxima by suprema and use one machine for all times.

\begin{theorem}[Refreshing experiment memory law]\label{thm:refresh}
For every fixed $d\ge1$, constants $0<c_d\le C_d<\infty$ exist such that,
for every \eqref{eq:refresh-range}, every integer $M\ge1$, every
$N\in\N\cup\{\infty\}$, and every $1\le t\le N$,
\begin{align}
 c_d\epsilon^2M^{-2/d}
 &\le R_{M,t}^{\rm av}\le R_{M,t}^{\rm max}
 \le C_d\epsilon^2M^{-2/d},\label{eq:refresh-checkpoint}\\
 c_d\epsilon^2M^{-2/d}
 &\le\mathfrak R_{M,N}^{\rm av}
 \le\mathfrak R_{M,N}^{\rm max}
 \le C_d\epsilon^2M^{-2/d}.
 \label{eq:refresh-causal}
\end{align}
The upper bounds are realized by one deterministic finite-label filter
using a fixed codebook for all positive times. The lower bounds allow
independent public and private coding randomization. At $\epsilon=0$ all
four risks are zero with one label. Constants are uniform in time, horizon,
calibration within \eqref{eq:refresh-range}, and the integer budget.
\end{theorem}

The proof will derive the cover, the actual attained density, and the
stable update separately. This also supplies a concrete dynamic realization
of the general emulation theorem.

\subsection{Explicit posterior recursion}
Set $p_0=\nu$. Given the previous posterior $p$, the distribution just before
the next report is
\begin{equation}\label{eq:refresh-q}
 q=(1-\gamma)p+\gamma\nu,
 \qquad q_i\ge m_*:=\frac{4}{5(d+1)}.
\end{equation}
For a command and report put
\[
 r_0=1,\quad r_i=1+y\epsilon\tanh u_i\ (i\ge1),
 \quad D_y=q_0+\sum_{i=1}^dq_ir_i.
\]
Then the probability of that report is $D_y/2$, and
\begin{equation}\label{eq:refresh-update}
 F(p,u,y)_0=\frac{q_0}{D_y},
 \qquad F(p,u,y)_i=\frac{q_ir_i}{D_y}\quad(i\ge1).
\end{equation}
Every denominator belongs to $[1-\epsilon,1+\epsilon]$. These formulas
hold for every belief in the full simplex and every admitted input.
Predictable actions reveal no further hidden-state information once their
past is conditioned upon, as in Proposition~\ref{prop:bayes}.

\begin{lemma}[Uniform contraction]\label{lem:refresh-contraction}
For any two simplex beliefs and the same $(u,y)$,
\begin{equation}\label{eq:refresh-contraction}
 \norm{F(p,u,y)-F(\tilde p,u,y)}_1
 \le\rho\norm{p-\tilde p}_1,
 \quad
 \rho=2\frac{1+\epsilon}{1-\epsilon}(1-\gamma)\le\frac23.
\end{equation}
\end{lemma}
\begin{proof}
For fixed positive likelihood vector $r$, write $B_r(q)=\diag(r)q/(q^Tr)$.
Subtract the two normalized vectors by adding and subtracting
$\diag(r)\tilde q/(q^Tr)$. The first term has norm at most
$(r_{\max}/r_{\min})\norm{q-\tilde q}_1$. The second has norm
$|q^Tr-\tilde q^Tr|/(q^Tr)$, at most the same bound. Thus
$\Lip_1(B_r)\le2r_{\max}/r_{\min}$.
Equation~\eqref{eq:refresh-q} multiplies differences by $1-\gamma$.
Use $r_{\max}\le1+\epsilon$, $r_{\min}\ge1-\epsilon$, and
\eqref{eq:refresh-range} to obtain the result.
\end{proof}

\begin{lemma}[An invariant shrinking domain]\label{lem:refresh-domain}
Let
\begin{equation}\label{eq:refresh-radius}
 a_\epsilon=\frac{2\epsilon}{\gamma(1-\epsilon)},
 \qquad
 D_\epsilon=\{p\in\Delta_d:\norm{p-\nu}_1\le a_\epsilon\}.
\end{equation}
Then $\nu\in D_\epsilon$, $F(D_\epsilon,u,y)\subset D_\epsilon$, and
$a_\epsilon\le(10/3)\epsilon$. Every attainable posterior belongs to
$D_\epsilon$, for every time and every admitted history.
\end{lemma}
\begin{proof}
For any pre-report belief $q$,
\[
 \norm{B_r(q)-q}_1
 =\sum_iq_i\frac{|r_i-q^Tr|}{q^Tr}
 \le\frac{2\epsilon}{1-\epsilon}.
\]
Consequently
\[
 \norm{F(p,u,y)-\nu}_1
 \le(1-\gamma)\norm{p-\nu}_1
          +\frac{2\epsilon}{1-\epsilon}.
\]
The right side is at most $a_\epsilon$ when $p\in D_\epsilon$, by its
definition. The numerical bound follows from
$\gamma\ge4/5$ and $1-\epsilon\ge3/4$. Start the induction at $\nu$.
\end{proof}

The set $D_\epsilon$ is a compact convex polytope for each fixed
calibration. Its representatives need not all be attainable at a particular
time; they are normalized beliefs with valid updates and belong to a proved
invariant domain. No extrapolation outside the probability simplex is used.
At $\epsilon=0$ the domain is the singleton $\{\nu\}$.

\subsection{The actual acquisition density}
Fix a preceding history, so its pre-report belief $q$ is fixed with
\eqref{eq:refresh-q}. For $\epsilon>0$ and a fixed sign $y$, the command
map
\[
 \Phi_{q,y}:u\longmapsto F(p,u,y)_+\in\R^d
\]
is one-to-one. Indeed its odds satisfy
\begin{equation}\label{eq:refresh-odds}
 \frac{p_i'}{p_0'}=\frac{q_i}{q_0}(1+y\epsilon\tanh u_i),
\end{equation}
so $u_i$ is recovered by applying $\operatorname{artanh}$ to
$\{p_i'q_0/(p_0'q_i)-1\}/(y\epsilon)$. The image of the command interior
has a smooth inverse. Boundary commands have zero probability under the
uniform acquisition law.

\begin{lemma}[Posterior Jacobian and unconditional density]\label{lem:refresh-density}
For $\epsilon>0$,
\begin{equation}\label{eq:refresh-jacobian}
 \left|\det D\Phi_{q,y}(u)\right|
 =\frac{q_0\prod_{i=1}^dq_i}{D_y^{d+1}}
       \epsilon^d\prod_{i=1}^d\sech^2u_i.
\end{equation}
Under independent uniform commands and actual reports, the law of
$p_{t,+}$ has a Lebesgue density bounded by $B_d\epsilon^{-d}$ at every
$t\ge1$, where one possible uniform constant is
\begin{equation}\label{eq:refresh-density-constant}
 B_d=2^{-d}(5/4)^{d+2}m_*^{-(d+1)}\cosh^{2d}(1).
\end{equation}
The bound is conditional on every preceding history as well as
unconditional.
\end{lemma}
\begin{proof}
Write $r_i'=y\epsilon\sech^2u_i$. In the $d$ nonzero-state coordinates,
\[
 D\Phi_{q,y}
 =\diag\left(\frac{q_ir_i'}{D_y}\right)
      -\left(\frac{q_ir_i}{D_y}\right)_i
       \left(\frac{q_jr_j'}{D_y}\right)_j^{\mathsf T}.
\]
The determinant lemma gives the determinant of the diagonal matrix times
\[
 1-\sum_{i=1}^d\frac{q_ir_i}{D_y}=\frac{q_0}{D_y}.
\]
Taking absolute values gives \eqref{eq:refresh-jacobian}.

Before any conditioning on the sign, the joint density of the uniform
command and that sign is
$2^{-d}D_y/2$. Change variables through the injective command map. Its
contribution to the posterior density is
\[
 \frac{2^{-(d+1)}D_y^{d+2}}
 {q_0\prod_{i=1}^dq_i\,\epsilon^d\prod_{i=1}^d\sech^2u_i}
 \le 2^{-(d+1)}(5/4)^{d+2}
        m_*^{-(d+1)}\cosh^{2d}(1)\epsilon^{-d}.
\]
Sum the two signs, allowing their images to overlap, to obtain
\eqref{eq:refresh-density-constant}. Finally integrate the conditional
bounds against the actual preceding-history law. The same bound survives
because it is independent of that history. Thus no event probability or
complementary command volume is discarded.
\end{proof}

This is a full-density upper bound for the actual law, which is sufficient
for a quantization lower bound. It is stronger than merely exhibiting a
nonzero differential at one command, but it does not assert that every
point of the invariant polytope is attained. Those are different statements.

\subsection{A fixed codebook and the proof of the theorem}
Project $D_\epsilon$ to its $d$ coordinates $p_+$. It lies in the cube
$\nu_++[-a_\epsilon,a_\epsilon]^d$. Partition this cube into
$k^d$ cells, $k=\lfloor M^{1/d}\rfloor$, and select one belief in
$D_\epsilon$ from each occupied cell. Fix those representatives once and
for all, along with a first-minimum tie convention. Empty cells need no
label. The discarded coordinate satisfies $p_0=1-\sum_{i=1}^dp_i$.
Hence within a cell the full-simplex $\ell^1$ distance is at most twice
the $\ell^1$ distance of its $d$ coordinates. The resulting codebook has
at most $M$ points and covering radius
\begin{equation}\label{eq:refresh-cover}
 r_M\le\frac{4da_\epsilon}{k}
 \le\frac{80d}{3}\epsilon M^{-1/d}.
\end{equation}
Choose the nearest representative in full $\ell^1$ distance. This preserves
the covering bound and is Borel. The initial state is exactly $\nu$;
thereafter the same codebook and the same update apply at every step.

\begin{proof}[Proof of Theorem~\ref{thm:refresh}]
For $\epsilon>0$, use \eqref{eq:refresh-update} at the retained belief and
then quantize with \eqref{eq:refresh-cover}. Lemmas~\ref{lem:refresh-contraction}
and \ref{lem:refresh-domain} justify the whole procedure on every history.
By Corollary~\ref{cor:contractive},
\begin{equation}\label{eq:refresh-error}
 \sup_{t\ge1}\norm{p_t-\hat p_t}_1
 \le3r_M\le80d\epsilon M^{-1/d}.
\end{equation}
Equation~\eqref{eq:refresh-distortion} bounds the excess by
$(\tau^2/d)(80d)^2\epsilon^2M^{-2/d}$. This proves a pathwise upper bound
for one machine, hence all the required mean and worst-history upper bounds.
A checkpoint encoder can use the same final nearest-centre procedure.

For the lower bound, Lemma~\ref{lem:refresh-density} and
Lemma~\ref{lem:small-ball} give
\[
 \inf_{c_1,\ldots,c_M}\E\min_j\norm{p_{t,+}-c_j}_2^2
 \ge \tfrac12(2B_dv_d)^{-2/d}\epsilon^2M^{-2/d}.
\]
Multiply by $\tau^2/d$ and use Theorem~\ref{thm:quantization}. The bound
holds for every $t\ge1$ under the one fixed exploration law, with coding
randomness allowed. Every causal encoder is a checkpoint encoder at each
time, proving the causal lower bound for finite or infinite horizons.
The inequality between mean and worst-history risks is immediate.
When $\epsilon=0$, both reports have probability $1/2$ independent of the
state and command. Refreshing preserves the uniform preparation, so
$p_t=\nu$ on every history. A single label predicts the terminal probe
exactly relative to the full-history oracle. This proves the zero case and
completes the uniform statement.
\end{proof}

\begin{corollary}[Persistent bits and double attenuation]\label{cor:double-attenuation}
For $b$ persistent bits, take $M=2^b$ in
\eqref{eq:refresh-causal}. With the fixed terminal probe the squared
excess has order $\epsilon^2 2^{-2b/d}$. If the actual terminal probe is
instead
\[
 \Pp(Z=1\mid X_t=i,J=j)=1/2+\tau\epsilon\1_{\{i=j\}},
\]
the order is $\epsilon^4 2^{-2b/d}$, with the same uniformity and the
same filter. At $\epsilon=0$, zero bits suffice in both models.
\end{corollary}
\begin{proof}
The first assertion substitutes the number of labels. In the second
experiment the physical readout difference is multiplied by $\epsilon$,
while the acquisition process, posterior law, and filter are unchanged.
Its squared excess is therefore multiplied by $\epsilon^2$. For the lower
bound the same scaling applies to arbitrary decoder centres after clipping
to the physical mean interval. The zero case was established above.
\end{proof}

\subsection{Feedback-uniform emulation in this model}
For fixed command $u$, the probability of the positive report under belief
$p$ is
\[
 g_+(p,u)=\tfrac12\left(1+
   \epsilon\sum_{i=1}^d((1-\gamma)p_i+\gamma\nu_i)\tanh u_i\right).
\]
For binary laws their total variation is the difference of their positive
probabilities. Thus
\begin{equation}\label{eq:refresh-report-lip}
 \norm{g(\cdot\mid p,u)-g(\cdot\mid\tilde p,u)}_{\TV}
 \le\tfrac12\epsilon(1-\gamma)\norm{p-\tilde p}_1.
\end{equation}
By Theorem~\ref{thm:emulation}, the actual report-path law and the
finite-state reference report-path law obey
\begin{equation}\label{eq:refresh-path-TV}
 \norm{P_N^\alpha-\widehat P_N^\alpha}_{\TV}
 \le\min\{1,40dN\epsilon^2(1-\gamma)M^{-1/d}\}
\end{equation}
for every common feedback controller. This statement is uniform over
controllers but is finite-horizon; it does not identify infinite path laws.
If a terminal probe is included in the compared observed path, its
additional conditional error is at most
$\tau\norm{p_t-\hat p_t}_1$, averaged over its random index. This term
is to be added to \eqref{eq:refresh-path-TV}. Prediction regret for that
probe instead uses its squared error, as in
\eqref{eq:refresh-distortion}; the two claims are different.

The bound vanishes for the acquisition report law when $\gamma=1$:
each new state is independently uniform, so previous belief errors have
no effect on the next acquisition report. The current posterior can still
have nontrivial geometry and require finite prediction memory for the
immediate terminal probe. This limiting case is a useful check on the
meaning of each estimate.

\subsection{Calibration and finite precision}
The theorem is uniform over \eqref{eq:refresh-range}, but the codebook and
update may use the known values of $\gamma$ and $\epsilon$. It is not a
calibration-blind theorem. Numerically approximating each belief update by
an admissible belief within $\eta$ in $\ell^1$ changes
\eqref{eq:refresh-error} to $3(r_M+\eta)$. The physical excess is at most
$9\tau^2(r_M+\eta)^2/d$. Thus choosing $\eta$ proportional to
$\epsilon M^{-1/d}$ preserves the order. This is an accuracy requirement;
it does not assert zero workspace or a universal bit-complexity bound for
computing the update.

For rational calibration one can specify the invariant polytope by rational
linear inequalities. Representatives can be selected in its occupied grid
cells by a finite feasibility procedure; centres and a membership/tie rule
form a finite read-only program. With nonrational calibration the theorem
uses the declared calibrated real-number model; a numerical implementation
must supply its own certified approximation. The only persistent acquired
information after a report is its representative label.

Finally, the refresh lower bound and the fixed positive state floor are
substantive. If $\gamma$ tends to zero, the constant $2/3$ need not remain
valid and the argument does not give a horizon-uniform theorem. That case
is not hidden in the calibration interval.
